\subsection{Method-structure projection}
\label{sec:method-structure-projection}

Scientific meaning coordinates describe what a source claims. Structural
learning downstream of that also needs a source-local account of \emph{how the
source's method transforms its problem}, so the representation is extended
with a versioned method-structure projection. It is derived from a
provenance-bound method-realization record of the kind produced by
source-local epistemic reconstruction~\cite{orion2025paper1}, and retains
exact source, document, version and span identity, the target role and initial
obstruction, assumptions and preconditions, representation changes and
auxiliary objects where the source supports them, the transformation graph,
invariants, progress and termination semantics, reconstruction to the original
target, failure and negative-result statements, explicit author rationale when
the source actually supplies it, and coordinates typed as unknown otherwise.

This projection stays source-local. In particular, the absence of an author
rationale is not permission to invent one. The projection stores the exact
realization and span digests that licensed each structured interpretation, and
carries no authority to declare that two methods belong to the same formal
family.

\subsection{Cross-domain method alignment and plural views}
\label{sec:method-alignment}

Given two source-local projections, a claim-relative mapping $\mu$ may be
proposed by comparing declared coordinates such as preconditions,
assumptions, invariants, progress, termination and reconstruction. The
alignment outcome is one of aligned, related, obstructed, or unresolved. The
record preserves both projection identities, source-local names and meanings,
unmatched assumptions, coordinate-level support and obstruction, and the exact
reasons for refusing a merge. An aligned outcome means only that the declared
source-local coordinates agree for the proposed mapping purpose; establishing
formal equivalence is outside this mechanism's authority.

A clean vocabulary-changing example is numerical bisection against monotone
threshold calibration. Their operation names differ, but under the narrow
purpose ``bounded monotone localization'' both require an ordered domain and a
bracketed monotone target, preserve a target-containing bracket, make progress
by shrinking the admissible interval, terminate at a tolerance, and
reconstruct a representative in the original target space. A candidate
alignment may therefore be recorded while each donor's vocabulary and lineage
is preserved.

A false-merge control uses the same-looking midpoint-and-discard-half surface
sequence on a non-monotone response. The monotonicity and bracketing contract
no longer holds and the invariant can fail, so the correct result is an
obstruction; topology or lexical similarity cannot override that coordinate. A
third control uses an opaque recovery procedure whose source does not
establish progress or reconstruction semantics. Even when its known effects
resemble a rollback method, the alignment stays unresolved rather than filling
the missing coordinates from analogy.

Plural views are therefore first-class outcomes: structurally suggestive
sources need not collapse into one global method. An obstruction certificate
records exactly which coordinates prevent a merge and keeps both source
projections recoverable.

\subsection{Downstream structural-learning boundary}

What this mechanism owns is source-local method projection and candidate
plural alignment. Claim-relative structural reduction, substitution and
method-family membership are outside it, as is any learned distribution over
these versioned objects and any generated method structure proposed from such
a distribution. Neither learned similarity nor generated structure changes a
source projection or confers scientific authority on it. The bounded pilot
reported in Section~\ref{sec:method-structure-eval} tests only representation
and alignment non-vacuity; it is not evidence of general extraction from
arbitrary papers.

\endinput
